\section{Supplementary Experimental Methods}
\label{sec:supp-methods}

\noindent\textbf{Build lineage and deployment.}
Current source-recovery code is committed as \texttt{d99ec92}. The
artifact records per-binary SHA-256 fingerprints and configuration for
each run. The current nine-service continuation, recovery, storage, and
short throughput series use this source-recovery rule. The initial
fourteen-service paired runs use the same single-obligation recovery
rule before the batch-result tag was separated from the direct-result
tag; the current cross-type codec regressions cover that separation.
Members in the current storage, in-memory continuation, and short
throughput controls run with GOGC=200 and GOMAXPROCS=16; this override is
not applied to committee, gateway, or wallet processes.

Archived operating-range runs record
base \texttt{e3870e1} and implementation \texttt{a083ff6}; original
continuation records \texttt{f0c0a274}; capital records base
\texttt{8e1fb40} and implementation \texttt{5807546}; organization paths
record base \texttt{bb65af69} and verified source \texttt{13be856}.

The archived batch integration uses \texttt{f82f8cd}; structural repair
counts were repeated on the current code. The source-failure checkpoint
\texttt{c1ebf2d} uses the superseded late-source rule and does not support
the current recovery claims.

Mac trials share one M4 Max host with 16 CPU cores and 64\,GiB RAM;
structural microbenchmarks run on Windows amd64. Current memory chains
and short throughput trials use in-memory service state and NoSync
wallets. Recovery trials use disk-backed committee application state
and BlockStore with in-memory members and gateway. Storage controls
change only member commits between bbolt NoSync and synchronous
transactions. Each run starts from fresh genesis.

Six additional fourteen-service disk-backed chain runs use the same
source-recovery rule and user-paid FUEL. Their three fast 100-hop runs
take 668.960--688.080\,ms; the inter-hop confirmation runs take
114.496--118.742\,s. They are a separate storage/deployment configuration
from the nine-service results in the main paper.

\noindent\textbf{Timing and aggregation.}
Certified-receipt latency generally runs from the first actual wallet
HTTP send through recipient verification and local atomic recording.
Capital and source-failure drivers include a small amount of pre-send
construction; their comparisons remain within the same driver family.
A local NoSync transaction completion is not a synchronous durability
receipt. Public and member completion timestamps record observation of
authenticated execution, rather than the committee's internal commit.
The sender's single monotonic clock measures cross-site ACK latency,
including the return path. A recipient can start its successor before
the preceding sender receives that ACK. Offered targets, actual send
rates, and completion rates including drain are retained separately.
Captions distinguish
transaction quantiles from statistics across runs. Within-run sample
variability is not an estimate of cross-run uncertainty.

\noindent\textbf{Externally funded authorization controller.}
The capital experiments sample member availability every 200\,ms.
Let $\rho$ be the declared CAL authorization demand per second, counting
both parent and child requests in each workload unit, and let $u$ count
unique requests whose certification remains unfinished. The controller
sets
\begin{align*}
 b &= 200+100u,\\
 H_{\mathrm{low}} &= \rho\,(0.8\,\mathrm{s})+b,\\
 H_{\mathrm{target}} &= 1.25H_{\mathrm{low}}.
\end{align*}
Here the watermarks and burst allowance are in CAL. When the minimum
observed member availability $A_{\min}$ falls below $H_{\mathrm{low}}$,
the proposed grant increment is
\[
 \Delta B=100\left\lceil
 \frac{1.5(H_{\mathrm{target}}-A_{\min})}{100}
 \right\rceil,
\]
subject to the configured authorization cap and available external
funding. At most one increment is in flight. Its public command debits
the external source and credits the reserve and grant atomically;
following members add the resulting share difference without clearing
existing approvals. The experiment supplies $\rho$ in advance rather
than estimating unknown production demand. Controller decisions are
requests for actual funded transfers, not a mechanism for creating
backing capital.

\noindent\textbf{Capital workload accounting.}
Each condition has one 300-s run with one Worker per member and sufficient
user FUEL. Starting from 3,600 CAL, the constant, rate-step, and delayed-parent
conditions end at 14,400, 21,600, and 23,300 CAL after external transfers
of 10,800, 18,000, and 19,700 CAL. Each completes 6,000 two-hop units,
or 12,000 payments, spending 1,128,000 user FUEL. Corresponding fixed-high
controls complete; the fixed-low condition closes 46 of 6,000 units,
leaves 5,826 unstarted, and records 64 partially approved payments.
These categories have different units and are not summed as one
payment count. The per-run tables preserve the full seven-condition
comparison without implying three repetitions per condition.

\noindent\textbf{Identities, resource use, and retained history.}
The multi-identity configurations use shared clients; 2,048 wallet
identities do not represent 2,048 separate operating system processes. Comparing
independent inputs with ten-hop payments, the three-run medians are
12.67 versus 15.24 core-ms per payment and 59.32 versus 71.83\,KiB of
HTTP payload per payment. Memory reporting separates total node RSS
and retained history from active obligations, pending work, and outboxes.
An increasing historical footprint is therefore not reported as an
increasing payment backlog. Per-run resource summaries and their archived source locations are included in the artifact index.

\noindent\textbf{Endpoint-specific Lightning reference.}
The archived deployment uses LND v0.21.3-beta and Bitcoin Core 31.1 in regtest, with synchronous LND persistence and its 50-ms channel-commit setting. Each of three fresh deployments has two nodes and a direct channel. After 20 warm-up transfers, each run issues 100 serial, alternating 10,000-satoshi payments using prepared invoices,
with receiver \texttt{SETTLED} as the completion event; the reported
P50/P95 are 282.12/320.17\,ms. A separate Next run sends 300 payments
at 100 payments/s for 3\,s using in-memory service state and NoSync
wallet storage. Recipient-certified receipt has P50/P95 of
0.9625/1.2075\,ms. Workload concurrency, durability, and completion
semantics differ, so these records are deployment references rather
than a normalized acceleration ratio.


\begin{table}[t]
  \caption{Append-only compensation records and authorized representation
    updates (C3). Both close the account through the same economic
    transition; they differ in the stored historical representation.}
  \label{tab:c3-compare}
  \centering
  \small
  \begin{tabular}{@{}p{0.2\columnwidth}p{0.37\columnwidth}p{0.37\columnwidth}@{}}
    \toprule
    & Append-only record & Representation update (C3) \\
    \midrule
    Economic closure
      & Reserve debit and obligation state in the public command;
        successors unchanged.
      & Same transition and atomic step, committed with a revision task. \\
    Historical representation and original commitments
      & Original body unchanged; the funding source of a consumed output
        is found through an index of compensation records.
      & Stored body names the reserve debit; transaction identity, owner
        authorization, and complete block identity $(H_b,P_b)$ are
        unchanged and verified. \\
    Actual readers
      & Any reader of the command history and records.
      & Repair construction, installation, and local inspection read the
        revised body; wallets, synchronization, and replay read original
        commands. \\
    Storage and replay
      & Original blocks plus records; replay runs original commands and
        compensation events.
      & Original blocks plus revised bodies, revisions, and tasks; replay
        runs original commands and repair events; same-block repairs
        share part adaptation and installation advances to the latest
        committed revision. \\
    Progress dependency
      & Public execution of the compensation record.
      & Compensation commits with its validated revision and requires
        three adaptation contributions. \\
    Status
      & Design baseline; not implemented as a mode and not measured.
      & Implemented; costs in the main paper. \\
    \bottomrule
  \end{tabular}
\end{table}

\subsection{Archived Operating-Range Series}
The archived operating-range series complements the
continuation results. One active organization uses sponsored fees,
reused addresses, and independent final inputs. Five runs of 150,000
payments each all finish (Fig.~\ref{fig:operating}). At a target of
2,000 payments/s, two runs achieve completion rates of
1,948.66--1,950.45 payments/s, including drain, with receipt P50s of
2.521--2.566\,ms. Raising the target to 2,200 yields
2,103.63--2,106.29 completed payments/s. At 2,400, the achieved rate is
2,174.68, P50/P95 rise to 7.174/89.38\,ms, and the 4,096-entry pending
admission limit records 1,635 hits. The operating curve pairs achieved completion rates with their corresponding queue pressure.




\begin{figure}[ht]
  \centering
  \includegraphics[width=0.62\columnwidth]{operating-point.pdf}
  \caption{Operating points versus target offered load. Top: completion
    throughput (payments/s), including drain. Bottom: receipt latency (ms).
    Each point is one 150,000-payment run:
    two at target 2,000, two at 2,200, and one at 2,400.
    P50/P95 summarize transactions within each run, not cross-run
    uncertainty. All services share one host.}
  \label{fig:operating}
\end{figure}
